\subsection*{9. An appraisal of teaching activity}
The ninth section returns on p.~49 to the university teaching begun in
1970 and to the change of professional milieu described at the beginning
of the report. Grothendieck recognizes what this experience taught him
and its part in renewing his mathematical outlook, then calls his
research-level teaching a failure by the criterion of creating a place
where something happens and something ``passe'' among people engaged in
a common adventure. The two exceptions he records, around 1974--76, are
Ladegaillerie's isotopy work already discussed in \S3, and
Contou-Carrère's work on local and global Jacobians over general base
schemes. The rest of the appraisal concerns teaching the rudiments of
the mathematician's craft to some twenty research students. Its final
paragraph on p.~50 qualifies the negative account by recalling work
better than merely acceptable, sometimes pursued beyond a single
project and carried through with taste. He nevertheless anticipates few
opportunities for these students to practice that craft professionally.
This reassessment develops the disappointment described at the end of
\S8 and supplies the immediate background to his application for
secondment in \S10; it is the author's account at the date of the report.
(\href{https://grothendieck.umontpellier.fr/web/viewer.html?file=../119.pdf\#page=52}{Copy 1, pp.~49--50, beginning at site 51};
\href{https://grothendieck.umontpellier.fr/web/viewer.html?file=../119.pdf\#page=53}{closing paragraph on p.~50, site 52}.)
